\section{Description}

The task is to find all longest common substrings of two strings. To do this, a generalized
suffix tree is built for the string $S_1 \# S_2 \$$, where instead of ordinary separator
characters the implementation uses the special codes 256 and 257
(\texttt{FIRST\_SEPARATOR} and \texttt{SECOND\_SEPARATOR}). Therefore the service separators
cannot coincide with any character of the input strings.

The suffix tree is built by Ukkonen's online algorithm. The transitions of a node are stored in
\texttt{std::map}, so every step costs $O(\log \sigma)$ and the whole construction takes
$O(n \log \sigma)$, where $n = |S_1| + |S_2| + 2$ and $\sigma \le 258$ is the size of the alphabet.
After the construction
the tree is traversed. For every node it is determined whether its subtree contains suffixes
of the first string and suffixes of the second string. If both kinds of suffixes are present,
the string spelled by the path from the root to this node is a common substring.
The deepest such node gives the length of the longest common substring.

\pagebreak

\section{Source code}

\subsection{Tree structure}

The class \texttt{SuffixTree} stores its nodes in a vector. Every node (the private structure
\texttt{Node}) stores the transitions by symbols (\texttt{children}), the suffix link
(\texttt{suffixLink}) and the bounds of the incoming edge in the combined string
(\texttt{start}, \texttt{end}). For leaves the right end of the edge equals the special value
\texttt{LEAF\_END}; the actual end of a leaf edge is given by the current member
\texttt{leafEnd\_}. New nodes are created by \texttt{newNode}, the length of an edge is computed
by \texttt{edgeLength}, and \texttt{child} returns the transition by a symbol or \texttt{NONE}.

The alphabet of the tree consists of integers. Ordinary bytes of the input strings fall into
the range $[0,255]$, and the separators have the values 256 and 257. The combined text is
built by the function \texttt{buildCombinedText}.

\subsection{Ukkonen's algorithm}

When the next symbol is added by \texttt{extend}, the algorithm maintains the active node
\texttt{activeNode\_}, the active edge \texttt{activeEdge\_}, the active length
\texttt{activeLength\_} and the number of remaining suffixes \texttt{remaining\_}. If the required
transition is missing, a leaf is created. If the transition exists but the next symbol on the
edge differs, the edge is split and an internal node is created. Suffix links make it possible
not to start processing every suffix from scratch.

\subsection{Finding the answer}

The method \texttt{findLcs(separator)} returns the structure \texttt{LcsResult} with the length
of the answer (\texttt{length}) and the sorted list of substrings (\texttt{substrings}).
Instead of a recursive depth-first search, which overflowed the call stack on inputs such as
\texttt{aaaa\ldots}, the traversal is iterative and uses an explicit stack. First, the nodes are
visited in preorder; for every node its parent and its string depth are recorded. Then the nodes
are processed in reverse order, so every node is handled after all its descendants.

For every node a coverage mask is aggregated: whether its subtree contains suffixes of the first
string, of the second string, or of both strings. To restore the substring itself, the start
position of one suffix of the first string (\texttt{firstStart}) is propagated to the parent as
well. If an internal node is covered by both strings, its depth is compared with the current
maximum, and the corresponding string, built by \texttt{makeString}, is added to a
\texttt{std::set}. The set removes duplicates automatically and sorts the answers
lexicographically.

The function \texttt{longestCommonSubstrings} builds the tree for two non-empty strings and calls
\texttt{findLcs}; if one of the strings is empty, the answer has length 0. The function
\texttt{runLcs} reads two lines (a missing second line is treated as an empty string, so the
program prints 0) and writes the answer. The \texttt{main} function only calls \texttt{runLcs}.

\lstinputlisting[caption={include/suffix\_tree.hpp}]{../include/suffix_tree.hpp}

\lstinputlisting[caption={src/suffix\_tree.cpp}]{../src/suffix_tree.cpp}

\lstinputlisting[caption={include/lcs.hpp}]{../include/lcs.hpp}

\lstinputlisting[caption={src/lcs.cpp}]{../src/lcs.cpp}

\needspace{18\baselineskip}
\lstinputlisting[caption={app/main.cpp}]{../app/main.cpp}

\pagebreak

\section{Console}

\begin{alltt}
# make build PRESET=release
# build/release/lab5/lab5_main
xabay
xabcbay
3
bay
xab
\end{alltt}

\pagebreak
